% PAGES: 105-106

\medskip
\noindent\rule{\linewidth}{0.4pt}
\medskip

\noindent\textbf{\textit{One Period Binomial Model Example:}}\par
In a one period binomial model for the evolution of the price of an asset,
it is assumed that the price at time $\tau > t_0$ of an asset with price
$S_0$ at time $t_0$ will be either $uS_0$ or $dS_0$, where $d < u$.

This corresponds to a one period market model with two securities, i.e.,
the asset described above and cash, and two states, i.e., state
$\omega^1$, when the price of the asset at time $\tau$ is $uS_0$, and
state $\omega^2$, when the price of the asset at time $\tau$ is $dS_0$.

Assume, without restricting the generality of the model, that the cash
position at time $t_0$ is \$1. If $r$ denotes the continuously compounded
annualized risk--free rate of return of a deposit made at time $t_0$ and
maturing at time $\tau$, then the future value at time $\tau$ of \$1 at
time $t_0$ is $e^{r(\tau-t_0)} = e^{r\delta t}$, where $\delta t = \tau -
t_0$.

The price vector of the securities at time $t_0$ is

\begin{equation}
S_{t_0} \;=\; \begin{pmatrix} 1 \\ S_0 \end{pmatrix}
\end{equation}

and the payoff matrix at time $\tau$ corresponding to this one period
market model is

\begin{equation}
M_\tau \;=\; \begin{pmatrix} e^{r\delta t} & e^{r\delta t} \\ uS_0 & dS_0
\end{pmatrix}.
\end{equation}

If, as it was assumed, $d < u$, then the two securities from the one
period binomial model are nonredundant.

To see this, note that the securities are nonredundant if and only if
$\rank(M_\tau) = 2$. Note that $\rank(M_\tau) = 1$ if and only if the rows
of the matrix $M_\tau$ are scalar multiples of each other, i.e., if and
only if there exists a constant $c$ such that

\[
\begin{aligned}
\begin{pmatrix} uS_0 & dS_0 \end{pmatrix} \;=\; c\begin{pmatrix}
e^{r\delta t} & e^{r\delta t} \end{pmatrix}
&\iff
\begin{cases} uS_0 \;=\; ce^{r\delta t} \\ dS_0 \;=\; ce^{r\delta t}
\end{cases} \\[6pt]
&\iff u \;=\; d \;=\; \dfrac{ce^{r\delta t}}{S_0},
\end{aligned}
\]

which contradicts the assumption that $d < u$.

We conclude that $\rank(M_\tau) = 2$ and therefore that the two assets are
non--redundant. \hfill$\square$

\medskip
\noindent\rule{\linewidth}{0.4pt}
\medskip

Consider a portfolio made of the $m$ securities and consisting of
$\Theta_j$ units of asset $j$ at time $t_0$, $j = 1:m$. The value $V_{t_0}$
of the portfolio at time $t_0$ is

\begin{equation}
V_{t_0} \;=\; \sum_{j=1}^m \Theta_j S_{jt_0} \;=\; \Theta^t S_{t_0},
\end{equation}

where $\Theta = (\Theta_j)_{j=1:m}$ denotes the $m \times 1$ column vector
of securities positions; cf.\ (1.5) and since $S_{t_0} =
(S_{jt_0})_{j=1:m}$.

\begin{lemma}
Consider a one period market model with $m$ securities and $n$ market
states at time $\tau > t_0$ and with payoff matrix $M_\tau =
\mathrm{row}\,(S_{j\tau})_{j=1:m}$. Let $\Theta$ be the positions vector at
time $t_0$ of a portfolio made of the $m$ securities and assume that the
asset positions in the portfolio remain unchanged until time $\tau$. If
$V_\tau$ is the $1 \times n$ row vector of the possible values of the
portfolio at time $\tau$, then

\begin{equation}
V_\tau \;=\; \Theta^t M_\tau \;=\; \sum_{j=1}^m \Theta_j S_{j\tau}.
\end{equation}
\end{lemma}

\begin{proof}
Let $V_\tau^{\omega^k}$ be the value at time $\tau$ of the portfolio if
state $\omega^k$ occurs, for $k = 1:n$. Since the asset positions at time
$\tau$ are $\Theta_j$, $j = 1:m$, it follows that

\begin{equation}
V_\tau^{\omega^k} \;=\; \sum_{j=1}^m \Theta_j S_{j\tau}^{\omega^k} \;=\;
\Theta^t S_\tau^{\omega^k},
\end{equation}

where $S_{j\tau}^{\omega^k}$ is the value of asset $j$ at time $\tau$ if
state $\omega^k$ occurs, for $k = 1:n$, and $S_\tau^{\omega^k} =
\left(S_{j\tau}^{\omega^k}\right)_{j=1:m}$; see (3.1). Then, from (3.8), we
obtain that

\[
\begin{aligned}
V_\tau &= \left(V_\tau^{\omega^1}\ \ V_\tau^{\omega^2}\ \ \ldots\ \
V_\tau^{\omega^n}\right) \;=\; \left(\Theta^t S_\tau^{\omega^1}\ \ \Theta^t
S_\tau^{\omega^2}\ \ \ldots\ \ \Theta^t S_\tau^{\omega^n}\right) \\
&= \Theta^t\,\mathrm{col}\left(S_\tau^{\omega^k}\right)_{k=1:n} \\
&= \Theta^t M_\tau,
\end{aligned}
\]

since $M_\tau = \mathrm{col}\left(S_\tau^{\omega^k}\right)_{k=1:n}$; cf.\
(3.2).

Using (1.9) and the row form $M_\tau = \mathrm{row}\,(S_{j\tau})_{j=1:m}$
of the payoff matrix $M_\tau$, see (3.3), we conclude that

\[
V_\tau \;=\; \Theta^t M_\tau \;=\; \sum_{j=1}^m \Theta_j S_{j\tau}.
\]
\end{proof}

\begin{definition}
A derivative security is replicable in a market model with $m$ securities
and $n$ market states at time $\tau > t_0$ if there exists a portfolio
made of the $m$ securities which has the same value as the derivative
security in every state of the market at time $\tau$.
\end{definition}

\begin{lemma}
Consider a market model with $m$ securities and $n$ market states at time
$\tau > t_0$ with payoff matrix $M_\tau = \mathrm{row}\,(S_{j\tau})_{j=1:m}$.
Let $s_\tau$ be the price vector at time $\tau$ of a replicable derivative
security, and let $\Theta = (\Theta_j)_{j=1:m}$ be the positions vector of
the replicating portfolio. Then,

\begin{equation}
s_\tau \;=\; \Theta^t M_\tau \;=\; \sum_{j=1}^m \Theta_j S_{j\tau}.
\end{equation}
\end{lemma}

We note that, if a derivative security is replicable, its value at time
$t_0$ is not necessarily uniquely determined, unless the market is
arbitrage--free; see Theorem 3.1.

\begin{proof}
By Definition 3.2, if $V_\tau$ is the price vector at time $\tau$ of the
replicating portfolio of a security with price vector $s_\tau$ at time
$\tau$, then $V_\tau = s_\tau$. From (3.7), we find that

\[
s_\tau \;=\; V_\tau \;=\; \Theta^t M_\tau \;=\; \sum_{j=1}^m \Theta_j
S_{j\tau}.
\]
\end{proof}
